\documentclass{article}
\usepackage{cancel}
\usepackage{gensymb}
\usepackage{tikz}
\usepackage{amsmath}
\usepackage{ amssymb }
\usepackage{listings}
\usepackage{geometry}
\usepackage{graphicx}
\usepackage{amsmath}
\usepackage{amsthm}
\usepackage{empheq}
\usepackage{mdframed}
\usepackage{booktabs}
\usepackage{lipsum}
\usepackage{color}
\usepackage{psfrag}
\usepackage{pgfplots}
\usepackage{bm}
\usepackage{xcolor}
\usepackage{minted}       
\usepackage{verbatim}    
\usepackage{siunitx}
\usepackage{fancyhdr}

\geometry{a4paper}



\usetikzlibrary{angles,quotes}
\usetikzlibrary{arrows.meta,calc}

% --- visual style for the cell ---
\newmdenv[
  backgroundcolor=gray!5,
  linecolor=gray!50,
  roundcorner=5pt,
  innerleftmargin=10pt,
  innerrightmargin=10pt,
  innertopmargin=10pt,
  innerbottommargin=10pt
]{codecell}

\newif\ifcodefast
\codefasttrue % change whether to write code fancy or just plain

\ifcodefast
  % ---------- FAST MODE ----------
  \newenvironment{codeblockPython}
    {\VerbatimEnvironment
     \begin{codecell}
     \begin{Verbatim}}
    {\end{Verbatim}
     \end{codecell}}
\else
  % ---------- FULL MINTED MODE ----------
  \newenvironment{codeblockPython}
    {\VerbatimEnvironment
     \begin{codecell}
     \begin{minted}{python3}}
    {\end{minted}
     \end{codecell}}
\fi

\geometry{a4paper}
\renewcommand{\thesubsubsection}{\alph{subsubsection})}
\setlength{\parindent}{0pt}
\DeclareSIUnit{\Msun}{M_{\odot}}
\DeclareSIUnit{\Lsun}{L_{\odot}}
\DeclareSIUnit{\Mpc}{Mpc}


\def\hwNUM{13}


\title{Astro210 HW \hwNUM}
\author{Pierson Lipschultz}


\begin{document}
\maketitle
\setcounter{section}{\hwNUM}

\pagestyle{fancy}

\fancyhf{}
\renewcommand{\headrulewidth}{.0pt}  
% \fancyhead[L]{{\textbf{Pierson Lipschultz Astro 210 HW \hwNUM, {\today}}}}
\fancyfoot[C]{\thepage}


\subsection{}


\subsubsection{}
\hrule
\vspace{.5cm}

Given \dots 

\[E(t) = \frac{4\sigma_{SB}}{c}\left(T^4V\right), quad P(t) = \frac{4\sigma_{SB}}{3c}T^4\]

We find \dots
\[ \frac{dE}{dt}\frac{4\sigma_{SB}}{c}\left(T^4V\right) = \frac{4\sigma_{SB}}{c}\left(4T^3V\frac{dT}{dt} + T^4 \frac{dV}{dt}\right)\]

\[\boxed{\frac{4\sigma_{SB}}{c}\left(4T^3V\frac{dT}{dt} + T^4 \frac{dV}{dt}\right) = -\frac{4\sigma_{SB}}{3c}T^4 \frac{dV}{dt}}\]

\subsubsection{}
\hrule
\vspace{.5cm}

Given \dots 
\[\frac{4\sigma_{SB}}{c}\left(4T^3V\frac{dT}{dt} + T^4 \frac{dV}{dt}\right) = -\frac{4\sigma_{SB}}{3c}T^4 \frac{dV}{dt}\]

We find \dots
\[\cancel{\frac{4\sigma_{SB}}{c}}\left(4T^3V\frac{dT}{dt} + T^4 \frac{dV}{dt}\right) = -\cancel{\frac{4\sigma_{SB}}{c}}T^4 \frac{1}{3}\frac{dV}{dt}\]

\[\left(4\frac{\cancel{T^{3}}}{T^{\cancel{4}}}V\frac{dT}{dt} + \cancel{T^4} \frac{dV}{dt}\right) = -\cancel{T^4} \frac{1}{3}\frac{dV}{dt}\]


\[4\frac{1}{T}V\frac{dT}{dt} = -\frac{1}{3}\frac{dV}{dt} - \frac{dV}{dt}\]


\[\cancel{4}\frac{1}{T}V\frac{dT}{dt} = -\frac{\cancel{4}}{3}\frac{dV}{dt}\]

\[\frac{1}{T}V\frac{dT}{dt} = -\frac{1}{3}\frac{dV}{dt}\]

\[\boxed{\frac{1}{T}\frac{dT}{dt} = -\frac{1}{3V}\frac{dV}{dt}}\]

\subsubsection{}
\hrule
\vspace{.5cm}

Given \dots 
\[\frac{1}{T}\frac{dT}{dt} = -\frac{1}{3V}\frac{dV}{dt}, \quad V(t) \propto a(t)^3\]
We find \dots 

\[\frac{dV}{dt} = 3a^2 \frac{ da}{dt}\]

\[\frac{1}{T}\frac{dT}{dt} = -\frac{1}{a\cancel{^3}}\cancel{a^2}\frac{da}{dt}\]

\[\boxed{\frac{1}{T}\frac{dT}{dt} = -\frac{1}{a}\frac{da}{dt}}\]

\subsubsection{}
\hrule
\vspace{.5cm}

Given \dots 
\[{\frac{1}{T}\frac{dT}{dt} = -\frac{1}{a}\frac{da}{dt}}\]
We find \dots
\[\frac{1}{T} =\ln(T)dT, \quad \frac{1}{a} =\ln(a)da\]

\[\boxed{\frac{d}{dt} \ln(T) = \ln(a)\frac{d}{dt}}\]

\subsubsection{}
\hrule
\vspace{.5cm}

Given \dots 
\[\lambda = \frac{2900 \mu \text{m}}{T}\]

\[T = \frac{2900 \mu \text{m}}{\lambda}\]

We find \dots 

\[{\frac{d}{dt} \ln\left(\frac{2900 \mu \text{m}}{\lambda}\right) = \ln(a)\frac{d}{dt}}\]

\[{\frac{d}{dt} \ln({2900 \mu \text{m}}) - \ln({\lambda}) = \ln(a)\frac{d}{dt}}\]
\[\boxed{\ln(\lambda) \propto \ln(a)}\]

\subsubsection{}
\hrule
\vspace{.5cm}
Given \dots 
\[\frac{1}{2} \left(\frac{dr}{dt}\right)^3 = \frac{GM}{r(t)} + k\]
\[\left(\frac{\dot{a}}{a}^2\right) = \frac{8\pi Gr_0^2}{3}P(t) + \frac{2\kappa}{r_0^2} \frac{1}{a(t)^2}\]
\[M- \frac{4\pi}{3} p(t) r(t)^3\]
\[\frac{4G\pi r(t)^2 p(t)}{3} + \kappa = 1/2 \frac{dr}{dt}^3\]

\[r(t)= a(t)r_0\]
\[\frac{1}{2}r^2_0 \dot{a}^2 = \frac{4}{3} G\pi (a(t)r_0)^2 + \kappa =\frac{1}{2} \left(\frac{dr}{dt}\right)^2\]

\[\boxed{\left(\frac{\dot{a}}{a}\right)^2 = \frac{8\pi}{3}p(t) + \frac{2k}{r^2_0} \frac{1}{a(t)^2}}\]


\subsubsection{}
\hrule
\vspace{.5cm}
Given \dots 
\[\left(\frac{\dot{a}}{a}\right)^2 = \frac{8\pi}{3}p(t) + \frac{2k}{r^2_0} \frac{1}{a(t)^2}\]
\[\left(\frac{\dot{a}}{a}\right)^2 = H(t)^2\]
\[k = 0\]
\[H(t)^2 = \frac{8\pi G}{3c^2} v_c\]

\[\frac{3c^2H(t)^2}{8G\pi} = v_c\]

\subsubsection{}
\hrule
\vspace{.5cm}


\subsubsection{}
\hrule
\vspace{.5cm}


\subsection{Precision cosmological measurements}
\hrule
\vspace{.5cm}

\subsubsection{}
\hrule
\vspace{.5cm}

Given \dots 
\[H_{0_1} = 67 \text{kms}^{-1} \text{Mpc}^{-1}, \quad H_{0_2} = 74 \text{kms}^{-1} \text{Mpc}^{-1}\]
\[V_{rec} = d * H_0\]
We find \dots 

\[\boxed{V_{rec_1} = \SI{7.504e+06}{\m / \s}, \quad V_{rec_2} = \SI{8.288e+06}{\m / \s}}\]

\subsubsection{}
\hrule
\vspace{.5cm}
Given \dots 
\[c = \lambda \nu\]
\[\lambda = \frac{c}{\nu}\]
\[\nu = 22.23508 \text{ GHz}\]
We find \dots 
\[\boxed{\lambda = \SI{1.348e-02}{\m}}\]

\subsubsection{}
\hrule
\vspace{.5cm}
Given \dots 

\[V_{rec_1} = \SI{7.504e+06}{\m / \s}, \quad V_{rec_2} = \SI{8.288e+06}{\m / \s}\]
\[\lambda_{obsv} = \lambda_{rest}(1 +z)\]
\[z = \frac{v}{c}\]

We find \dots 
\[z_1 = \SI{2.503e-02}{}, \quad z_2 = \SI{2.765e-02}{}\]
\[\nu_1 = \SI{2.279e+10}{1/\s},\quad \nu_2 = \SI{2.285e+10}{1/\s}\]
\[\lambda_1 = \SI{1.315e-02}{\m}, \quad \lambda_2 = \SI{1.312e-02}{\m}\]
\[\boxed{\Delta \lambda = \SI{3.347e-05}{\m}}\]


\subsection{Matter-dark energy equality}
\hrule
\vspace{.5cm}

\subsubsection{}
\hrule
\vspace{.5cm}
Given \dots 
\[\Omega_{\lambda_0} = (1+z_{eq})^3\Omega_{m_0}\]
We find \dots 
\[\frac{\Omega_{\lambda_0}}{\Omega_{m_0}}= (1+z_{eq})^3\]
\[\left(\frac{\Omega_{\lambda_0}}{\Omega_{m_0}}\right)^{1/3} -1 = z_{eq}\]
\[\boxed{z_{eq} \approx 3.931 \times 10^{-1}}\]

\subsubsection{}
\hrule
\vspace{.5cm}

Given \dots 
\[\ddot{a} = H_0\left[\frac{\Omega_{r,0}}{a^2} + \frac{\Omega_{m,0}}{a}+ \Omega_{\lambda,0}a^2\right]^{1/2}\]
\[1 + z = \frac{1}{a} \Rightarrow a = \frac{1}{1+z}\]
We find\dots 
\[\boxed{\ddot{a} = H_0 (0.86734)}\]


\subsubsection{}
\hrule
\vspace{.5cm}
Given \dots 
\[\frac{H(t)^2}{H_0^2} = \frac{\Omega_{r,0}}{a(t)^4} + \frac{\Omega_{m,0}}{a(t)^3} + \Omega_{\Lambda}\]
We find \dots 

\[{H(t)} = {H_0^2}\left[\frac{\Omega_{m,0}}{a(t)^3} + \Omega_{\Lambda}\right]^{1/2}\]

\[\boxed{H(t) = (H_0)^2 1.2083}\]



\subsection{Curvature}
\hrule
\vspace{.5cm}

Given \dots 
\[\alpha + \beta + \gamma = \pi + \frac{\kappa A}{r^2_{c,0}}, \quad \text{By spherical} \Rightarrow \kappa = 1\]
\[\alpha + \beta + \gamma = \theta_{tot}\]
\[A = \frac{\sqrt{3}}{4} \ell^2, \quad \ell = 1\]
We find \dots 
\[\theta_{tot} = \pi + \frac{A}{r^2}\]
\[\boxed{\theta_{tot} = \SI{3.141592664e+00}{\radian}}\]
\[\boxed{\theta_{tot} - \pi = \SI{1.066806e-08}{\radian}}\]
\hrule


\begin{center}\footnote{spent an hour procrastinating the actual homework making this}
\begin{tikzpicture}
\begin{axis}[
  title={Spherical triangle on a unit sphere},
  axis lines = box,
  scale = 1.4,
  xlabel={$x$},
  ylabel={$y$},
  zlabel={$z$},
  colormap/cool,
  view={135}{25},
  xmin=-.3, xmax=1,
  ymin=-.3, ymax=1,
  zmin=-.3, zmax=1,
]

  % --- Unit sphere as a surface ---
  \addplot3[
    surf,
    samples=3,
    domain=0:360,      % longitude
    y domain=0:180,    % colatitude
  ]
  (
    {cos(x)*sin(y)},   % x
    {sin(x)*sin(y)},   % y
    {cos(y)}           % z
  );


  % --- Edge AB: equator from (1,0,0) to (0,1,0) ---
  %   \addplot3[
  %     thick,
  %     color=white,
  %   ]
  %   (
  %     {cos(x)},          % x
  %     {sin(x)},          % y
  %     {0}                % z
  %   );
  % default domain is 0:360, so restrict:
  \addplot3[
    thick,
    color=white,
    domain=0:90
  ]
  (
    {cos(x)},
    {sin(x)},
    {0}
  );

  % --- Edge AC: meridian at y=0, from equator to north pole ---
  \addplot3[
    thick,
    color=white,
    domain=0:90
  ]
  (
    {cos(x)},          % x
    {0},               % y
    {sin(x)}           % z
  );

  % --- Edge BC: meridian at x=0, from equator to north pole ---
  \addplot3[
    thick,
    color=white,
    domain=0:90
  ]
  (
    {0},               % x
    {cos(x)},          % y
    {sin(x)}           % z
  );

  % --- Mark the vertices A, B, C ---
  \addplot3[
    only marks,
    mark=*,
    mark size=2pt,
    color=black,
  ]
  coordinates {
    (1,0,0)
    (0,1,0)
    (0,0,1)
  };

\end{axis}
\end{tikzpicture}
\end{center}





\subsection{Proper distances}
\hrule
\vspace{.5cm}
Given \dots 
\[r = \int_{t_{eq}}^{t_0} \frac{dt}{a(t)}\]
\[a(t) = \frac{t}{t_0}\]

We find \dots 
\[r = ct_0\int_{t_{eq}}^{t_0} \frac{dt}{t}\]

\[r = ct_0\ln(t) |_{t_{eq}}^{t_{t_0}}\]

\[r = ct_0(\ln(t_0) -\ln(t_{eq}))\]

\[r = ct_0\left[\ln\left(\frac{t_0}{t_{eq}}\right)\right]\]

\[r = ct_0\left[\ln\left(\frac{1}{a_{eq}}\right)\right]\]

\[r = ct_0\left(\ln\left(1+z\right)\right)\]

\[\boxed{r \approx \SI{4.329e+25}{\m} \approx \SI{ 1.403e+03}{\Mpc}}\]


\end{document}